Thirdly, it seems well worth stating the causes to which the rapid and uninterrupted rise in the prosperity of the widows' fund over a period of more than forty years must be attributed. I find that this rise was especially favoured by the combined effect of two circumstances: one had not been counted on at all in 1794, while the other, at least, ought not to have been counted on.

The first cause is the rise in the rate of interest, which began only a few years after that date. As I have shown, it was then believed that only 3 per cent could safely be counted on. But already in 1799 the Curatorium remarked, in the rescript of 23 May cited above, that opportunities for safe investment at 4 per cent were by no means rare, and even offered its own help with this. Somewhat later, however, the rate of interest generally rose to 5 per cent and remained at that level for the greater part of the widows' fund's capital over a long series of years.

The second cause is that, over that long period, the number of pensions remained \textit{below} the expected mean (15); indeed, one must say \textit{considerably} below the expected mean if, following the more plausible estimate on p. [145], one takes this to be 17 (namely, \(\tfrac{4}{7}\times26\) widows, with an addition of \(\tfrac{1}{6}\) on account of the orphans). It has already been mentioned above on p. [144, line 7] that this occurrence should by no means be regarded as something very extraordinary. But it would have been no more extraordinary if exactly the opposite had occurred and, for example, only 20—25 years after that date of 1794, the mean had been considerably and persistently exceeded on the other side. In such a case the rapid increase in the fund's assets would have given way to a very slow one, perhaps even to a standstill, or indeed to the need to resort to the receipt clause if a sharp fall in the rate of interest had occurred at the same time.

A higher number of widows has in fact already been present for a number of years, p. [126, line 23] and
\SourceLeaf{154}
p. [128, line 29]; but the rate of interest, though it has fallen far from its former level, is still high enough for the fund to remain able to meet that higher number of pensions. In general, this high number, even if it had risen by one or a couple of pensions more, is in itself, and insofar as it is regarded merely as an unusually large fluctuation above the mean of 15 or 17, far less serious than \textit{the} circumstance to which I now turn
\begin{quote}
that the mean number of pensions, whether 15 or 17 is regarded as the most plausible, is valid only so long as the association does not grow beyond its size in 1794, and that this validity has now ceased altogether, since its size, measured by the number of married members, is more than 60 per cent greater than at that date.
\end{quote}
\noindent\hspace*{\parindent}This measure cannot be applied to the interval between 1794 and 1845, since the number of married members is not known. But so far as can be inferred from changes in the total number of all members, the size of the association, taken as a whole and apart from some fluctuations back and forth, was no greater up to about 1831, but rather somewhat smaller, than in 1794; and only in the last few years has the increase become \textit{very considerable}. The association's increased size is therefore in no way causally connected with the recent high number of widows, as already remarked above [p. 128].

The fourth point still to be considered here, and, as it were, the keystone of the first division, is to make an estimate of the future course of the number of widows in general, insofar as such an estimate can be reached along the path followed so far. The total number of professors is now 57; of these, 51 participate in the widows' fund, and 42 of them are married. If one first makes the estimate on the hypothesis that seems most appropriate to the circumstances of the professors' widows' fund here, namely, that for every 7 existing marriages, 4 widows are to be reckoned, the calculation gives 24 widows. This means that, once the society of that size has reached the steady state, the average number of widows will be 24. According to \textsc{Kritter}'s opinion, a further sixth would have to be added on account of the orphans' pensions; I shall, however, allow only an eighth, thus setting the average number of pensions in the steady state \(=27\). Admittedly, the steady state is to be expected only after 45—50 years; but one would be greatly mistaken in regarding its being so far off as a strong reason for reassurance. For long before that, the limiting state has been approached so closely that the difference no longer matters much. Compare the figures given above on p. [141] for the gradual rise in the number of widows: without claiming strict accuracy, they can nevertheless give some idea of the course of events. If one then considers that the society's former size had already given 17 pensions as the normal number (at the ratio \(7:4\)), and that what is at issue here is therefore only the gradual development of the increase estimated at 10, one will readily understand me when I must assert that after only about 25 years one must be prepared for 25 pensions as the \textit{mean}. To this must now be added the widening of the range owing to fluctuations, to which it is impossible to set sharp limits. But this much is certain: a fluctuation of 7 pensions above and below the mean must be included in the estimate as something by no means extraordinary, since earlier precedents establish that a fluctuation at least as large in proportion has actually occurred. The result of these considerations is therefore that one must expect to find the number of pensions around the year 1870 between 18 and 32, without being able to determine in advance \textit{where} within this wide range; that reaching either extreme must not be regarded as something very extraordinary; and finally, that these figures must later be increased a
\SourceLeaf{155}
little further, so that after a longer time the range must be designated by 20—34. That is all that the calculus of probabilities can teach, \textit{so long as the society is considered, as it were, only in the aggregate.} Some may perhaps think that this is little! I should think otherwise. It is very important to know the extent of the expectations one is entitled to have, and not to abandon oneself to an unfounded sense of security. Let us consider the two extremes. It is \textit{possible} that after 25 years the number of pensions has still not exceeded 18, and remains at that level, or only very little above it, for quite some time longer. The number of pensions may also first fall from its present level (19), and thus, from now until 1870, average appreciably below 18. In this way the fund can gather strength with which, when the tide turns only after a long time, it can withstand even the extreme cases on the other side. The worst thing would actually be if these favourable assumptions were realised to an even greater extent; I mean, if in the meantime the number of pensions were once again to fall below 16, so that under the existing statute the fund would be burdened with a further pension increase of 10, 20, perhaps even 30 Rthl. — If, conversely, as is equally possible, the number of pensions were after only 25 years to have risen to or near 32, a rough estimate suffices to show that, apart from quite extraordinary inflows, the fund will be unable to meet such demands, but will be compelled, presumably even earlier, to declare its inability to pay.

It goes without saying that the estimates become somewhat more favourable if, instead of the ratio \(7:4\), one assumes \(2:1\). One would then after 25 years have to expect a range of 16—28 pensions, and 18—30 after a still longer time.

Finally, I must remark that it is tacitly assumed here that the size of the association will remain at its present level from now on. If it increases further, one must be prepared for proportionately still larger numbers; if, on the other hand, it decreases again, some deduction from the above results will also be permissible. But by the nature of the matter, it seems unlikely that a \textit{considerable} decrease could occur except after a long time.

